\par\Needspace{6\baselineskip}\medskip
\hypertarget{native-218-article159}{}\pdfbookmark[2]{159.}{article-159}
\begin{center}159.\end{center}
For the \emph{third} example we choose the following observations of Ceres. The
first was made at Bremen by the distinguished \textsc{Olbers}, the second at Göttingen by the distinguished \textsc{Harding}, and the third at Lilienthal by the distinguished \textsc{Bessel}:
\begin{center}\begingroup\footnotesize\setlength{\tabcolsep}{3pt}\renewcommand{\arraystretch}{1.3}
\begin{tabular}{|c|c|c|}
\hline
\shortstack{Local mean time\\of observation} & Right asc. & North decl. \\
\hline
1805 Sept. 5. $13^{\mathrm h}$ $8^{\mathrm m}$ $54^{\mathrm s}$ & $95^0\,59'\,25''$ & $22^0\,21'\,25''$ \\
1806 Jan. 17. $10$ $58$ $51$ & $101\quad18\quad40{,}6$ & $30\quad21\quad22{,}3$ \\
1806 May 23. $10$ $23$ $53$ & $121\quad56\quad7$ & $28\quad2\quad45$ \\
\hline
\end{tabular}
\endgroup\end{center}

The preceding two examples have already amply illustrated the methods for taking
parallax and aberration into account when the distances from the Earth are regarded
as entirely unknown. In this third example we therefore dispense with the extra
work and take approximate distances from the scholarly correspondence of the
distinguished de \textsc{Zach} (Vol. XI p. 284) to correct the observations for the effects of parallax and aberration. The following table gives these distances
together with the reductions derived from them:
\begin{center}\begingroup\footnotesize\setlength{\tabcolsep}{3pt}\renewcommand{\arraystretch}{1.3}
\begin{tabular}{l|c|c|c}
Distance of Ceres from the Earth & $2{,}899$ & $1{,}638$ & $2{,}964$ \\
\shortstack[l]{time taken for light\\to reach the Earth} & $23^{\mathrm m}49^{\mathrm s}$ & $13^{\mathrm m}28^{\mathrm s}$ & $24^{\mathrm m}21^{\mathrm s}$ \\
reduced time of observation & $12^{\mathrm h}45^{\mathrm m}5^{\mathrm s}$ & $10^{\mathrm h}45^{\mathrm m}23^{\mathrm s}$ & $9^{\mathrm h}59^{\mathrm m}32^{\mathrm s}$ \\
sidereal time in degrees & $355^0\,55'$ & $97^0\,59'$ & $210^0\,41'$ \\
parallax in right ascension & $+1''{,}90$ & $+0''{,}22$ & $-1''{,}97$ \\
parallax in declination & $-2{,}08$ & $-1{,}90$ & $-2{,}04$ \\
\end{tabular}
\endgroup\end{center}

\SourceLeaf{219}Thus, after correction for parallax and aberration and reduction of the times
to the Paris meridian, the data of the problem are as follows:
\begin{center}\begingroup\footnotesize\setlength{\tabcolsep}{3pt}\renewcommand{\arraystretch}{1.3}
\begin{tabular}{|c|c|c|}
\hline
 & Right asc. & Declination \\
\hline
1805 Sept. 5. $12^{\mathrm h}$ $19^{\mathrm m}$ $14^{\mathrm s}$ & $95^0\,59'\,23''{,}10$ & $22^0\,21'\,27''{,}08$ \\
1806 Jan. 17. $10$ $15$ $2$ & $101\quad18\quad40{,}38$ & $30\quad21\quad24{,}20$ \\
1806 May 23. $9$ $33$ $18$ & $121\quad56\quad8{,}97$ & $28\quad2\quad47{,}04$ \\
\hline
\end{tabular}
\endgroup\end{center}

From these right ascensions and declinations we derived the longitudes and
latitudes, using the obliquities of the ecliptic $23^0\,27'\,55''{,}90$, $23^0\,27'\,54''{,}59$,
$23^0\,27'\,53''{,}27$. The longitudes were then corrected for nutation, which was respectively
$+17''{,}31$, $+17''{,}88$, $+18''{,}00$, and afterward reduced to the beginning of 1806 by applying the precession $+15''{,}98$, $-2''{,}39$, $-19''{,}68$. Finally, for the reduced times
we took the positions of the Sun from the tables. In the longitudes, nutation was
omitted, while precession was added just as for the longitudes of Ceres. The latitude
of the Sun was neglected altogether. This gave the following numbers for use in the calculation:
\begin{center}\begingroup\footnotesize\setlength{\tabcolsep}{3pt}\renewcommand{\arraystretch}{1.3}
\begin{tabular}{l|c|c|c}
Time 1805 Sept. & $5{,}51336$ & $139{,}42711$ & $265{,}39813$ \\
$a$, $a'$, $a''$ & $95^0\,32'\,18''{,}56$ & $99^0\,49'\,5''{,}87$ & $118^0\,5'\,28''{,}85$ \\
$\beta$, $\beta'$, $\beta''$ & $-0\quad59\quad34{,}06$ & $+7\quad16\quad36{,}80$ & $+7\quad38\quad49{,}39$ \\
$l$, $l'$, $l''$ & $342\quad54\quad56{,}00$ & $117\quad12\quad43{,}25$ & $241\quad58\quad50{,}71$ \\
$\log R$, $\log R'$, $\log R''$ & $0{,}003\,1514$ & $9{,}992\,9861$ & $0{,}005\,6974$ \\
\end{tabular}
\endgroup\end{center}

The preliminary calculations explained in Arts. 136---143 now give the following:
\begin{center}\begingroup\footnotesize\setlength{\tabcolsep}{3pt}\renewcommand{\arraystretch}{1.3}
\begin{tabular}{l|c|c|c}
$\gamma$, $\gamma'$, $\gamma''$ & $358^0\,55'\,28''{,}09$ & $156^0\,52'\,11''{,}49$ & $170^0\,48'\,44''{,}79$ \\
$\delta$, $\delta'$, $\delta''$ & $112\quad37\quad9{,}66$ & $18\quad48\quad39{,}81$ & $123\quad32\quad52{,}13$ \\
$A'D$, $AD'$, $AD''$ & $15\quad32\quad41{,}40$ & $252\quad42\quad19{,}14$ & $136\quad2\quad22{,}38$ \\
$A''D$, $A''D'$, $A'D''$ & $138\quad45\quad4{,}60$ & $6\quad26\quad41{,}10$ & $358\quad5\quad57{,}00$ \\
$\epsilon$, $\epsilon'$, $\epsilon''$ & $29\quad18\quad8{,}21$ & $170\quad32\quad59{,}08$ & $156\quad6\quad25{,}25$ \\
\end{tabular}
\endgroup\end{center}
\SourceLeaf{220}
\begin{center}\begingroup\footnotesize\setlength{\tabcolsep}{3pt}\renewcommand{\arraystretch}{1.3}
\begin{tabular}{ll}
$\sigma=8^0\,52'\,4''{,}05$ & \\
$\log a=0{,}184\,0193_n$, & $a=-1{,}527\,6340$\\
$\log b=0{,}004\,0987$ & \\
$\log\varkappa=0{,}161\,1012$ & \\
$\log\lambda=9{,}916\,4090_n$ &
\end{tabular}
\par\smallskip
\begin{tabular}{r@{\;=\;}l}
$\log c$ & $2{,}006\,6735$\\
$d$ & $117{,}50873$\\
$\log e$ & $0{,}856\,8244$\\
$\log\varkappa''$ & $9{,}977\,0819_n$\\
$\log\lambda''$ & $9{,}732\,0127_n$.
\end{tabular}
\endgroup\end{center}

The time interval between the first and second observations is
$133{,}91375$ days, and that between the second and third is $125{,}97102$ days. Hence $\log\theta = 0{,}335\,8520$,
$\log\theta'' = 0{,}362\,4066$, $\log\dfrac{\theta''}{\theta} = 0{,}026\,5546$, $\log\theta\theta'' = 0{,}698\,2586$. In the following table we give the principal quantities
for the first three hypotheses formed in succession:
\begin{center}\begingroup\footnotesize\setlength{\tabcolsep}{3pt}\renewcommand{\arraystretch}{1.3}
\begin{tabular}{l|c|c|c}
 & I & II & III \\
\hline
$\log P = x$ & $0{,}026\,5546$ & $0{,}025\,6968$ & $0{,}025\,6275$ \\
$\log Q = y$ & $0{,}698\,2586$ & $0{,}739\,0190$ & $0{,}748\,1055$ \\
$\omega + \sigma$ & $7^0\,15'\,13''{,}523$ & $7^0\,14'\,47''{,}199$ & $7^0\,14'\,45''{,}071$ \\
$\log Qc\sin\omega$ & $1{,}154\,6650_n$ & $1{,}197\,3880_n$ & $1{,}206\,6327_n$ \\
$z$ & $7^0\,3'\,59''{,}018$ & $7^0\,2'\,32''{,}938$ & $7^0\,2'\,16''{,}900$ \\
$\log r'$ & $0{,}411\,4726$ & $0{,}412\,9371$ & $0{,}413\,2107$ \\
$\zeta$ & $160^0\,10'\,46''{,}74$ & $160^0\,20'\,7''{,}82$ & $160^0\,22'\,9''{,}42$ \\
$\zeta''$ & $262\quad6\quad1{,}03$ & $262\quad12\quad18{,}26$ & $262\quad14\quad19{,}49$ \\
$\log r$ & $0{,}432\,3934$ & $0{,}429\,1773$ & $0{,}428\,4841$ \\
$\log r''$ & $0{,}409\,4712$ & $0{,}407\,1975$ & $0{,}406\,4697$ \\
$\tfrac12(u''+u)$ & $262^0\,55'\,23''{,}22$ & $262^0\,57'\,6''{,}83$ & $262^0\,57'\,31''{,}17$ \\
$\tfrac12(u''-u)$ & $273\quad28\quad50{,}95$ & $273\quad29\quad15{,}06$ & $273\quad29\quad19{,}56$ \\
$2f'$ & $62\quad34\quad28{,}40$ & $62\quad49\quad56{,}50$ & $62\quad53\quad57{,}06$ \\
$2f$ & $31\quad8\quad30{,}03$ & $31\quad15\quad59{,}09$ & $31\quad18\quad13{,}83$ \\
$2f''$ & $31\quad25\quad58{,}45$ & $31\quad33\quad57{,}32$ & $31\quad35\quad43{,}32$ \\
$\log \eta$ & $0{,}020\,2496$ & $0{,}020\,3158$ & $0{,}020\,3494$ \\
$\log \eta''$ & $0{,}021\,1074$ & $0{,}021\,2429$ & $0{,}021\,2751$ \\
$\log P'$ & $0{,}025\,6968$ & $0{,}025\,6275$ & $0{,}025\,6289$ \\
$\log Q'$ & $0{,}739\,0190$ & $0{,}748\,1055$ & $0{,}750\,2337$ \\
$X$ & $-0{,}000\,8578$ & $-0{,}000\,0693$ & $+0{,}000\,0014$ \\
$Y$ & $+0{,}040\,7604$ & $+0{,}009\,0865$ & $+0{,}002\,1282$ \\
\end{tabular}
\endgroup\end{center}
\SourceLeaf{221}

We now denote the three values of $X$ by $A$, $A'$, $A''$, and the three values of $Y$
by $B$, $B'$, $B''$. We let the quotients obtained by dividing $A'B''-A''B'$, $A''B-AB''$,
$AB'-A'B$ by their sum be $k$, $k'$, $k''$, respectively, so that $k+k'+k''=1$. Finally, we denote the values of $\log P'$ and $\log Q'$ in the third
hypothesis by $M$ and $N$ (these would be the new values of $x$, $y$ if the fourth
hypothesis were to be derived from the third in the same way as the third was
derived from the second). From the formulas of Art. 120 it follows readily that the corrected value of $x$ is
$=M-k(A'+A'')-k'A''$ and that the other corrected value is given by $y = N-k(B'+B'')-k'B''$.
Calculation gives the former $=0{,}025\,6331$ and the latter $=0{,}750\,9143$. On these corrected values we now base the \emph{fourth hypothesis}, whose principal quantities are as follows:
\begin{center}\begingroup\footnotesize\setlength{\tabcolsep}{3pt}\renewcommand{\arraystretch}{1.3}
\begin{tabular}{r|r}
$\omega + \sigma=7^0\,14'\,45''{,}247$ & $\log r''=0{,}406\,2033$\\
$\log Qc\sin\omega=1{,}209\,4284_n$ & $\tfrac12(u''+u)=262^0\,57'\,38''{,}78$\\
$z=7^0\,2'\,12''{,}736$ & $\tfrac12(u''-u)=273\quad29\quad20{,}73$\\
$\log r'=0{,}413\,2817$ & $2f'=62\quad55\quad16{,}64$\\
$\zeta=160^0\,22'\,45''{,}38$ & $2f=31\quad19\quad1{,}49$\\
$\zeta''=262\quad15\quad3{,}90$ & $2f''=31\quad36\quad15{,}20$\\
$\log r=0{,}428\,2792$ & \\
\end{tabular}
\endgroup\end{center}

Between $2f'$ and $2f+2f''$ there is a difference of $0''{,}05$, which we distribute by setting
$2f=31^0\,19'\,1''{,}47$, $2f''=31^0\,36'\,15''{,}17$. If the elements themselves are now determined from the two
outer positions, the following numbers result:
\begin{center}\begingroup\footnotesize\setlength{\tabcolsep}{3pt}\renewcommand{\arraystretch}{1.3}
\begin{tabular}{ll}
True anomaly at the first position\quad .\,.\,.\,.\,.\,.\,.\,.\,.\,.\,.\,.\,.\,.\,.\,. & $289^0\,7'\,39''{,}75$\\
True anomaly at the third position\quad .\,.\,.\,.\,.\,.\,.\,.\,.\,.\,.\,.\,.\,.\,.\,. & $352\quad2\quad56{,}39$\\
Mean anomaly at the first position\quad .\,.\,.\,.\,.\,.\,.\,.\,.\,.\,.\,.\,.\,.\,. & $297\quad41\quad35{,}65$\\
Mean anomaly at the third position\quad .\,.\,.\,.\,.\,.\,.\,.\,.\,.\,.\,.\,.\,.\,. & $353\quad15\quad22{,}49$\\
Mean daily sidereal motion\quad .\,.\,.\,.\,.\,.\,.\,.\,.\,.\,.\,.\,.\,.\,.\,.\,. & $769''{,}6750$\\
Mean anomaly at the beginning of 1806\quad .\,.\,.\,.\,.\,.\,.\,.\,.\,. & $322^0\,35'\,52''{,}51$\\
Angle $\varphi$\quad .\,.\,.\,.\,.\,.\,.\,.\,.\,.\,.\,.\,.\,.\,.\,.\,.\,.\,.\,.\,.\,.\,.\,.\,.\,.\,.\,.\,.\,.\,.\,.\,. & $4\quad37\quad57{,}78$\\
Logarithm of the semimajor axis\quad .\,.\,.\,.\,.\,.\,.\,.\,.\,.\,.\,.\,.\,.\,.\,.\,. & $0{,}442\,4661$.
\end{tabular}
\endgroup\end{center}

Computing the heliocentric position from these elements for the time of the middle observation, we find the mean anomaly $326^0\,19'\,25''{,}72$, the logarithm of the radius
vector $0{,}413\,2825$, and the true anomaly $320^0\,43'\,54''{,}87$. The latter should differ from the \SourceLeaf{222} true anomaly at the first position by $2f''$, or from the true anomaly at the
third position by $2f$, and should therefore be $=320^{0}\,43'\,54''{,}92$, just as the
logarithm of the radius vector should be $=0{,}413\,2817$. The difference of $0''{,}05$ in the true anomaly
and that of eight units in this logarithm may be regarded as insignificant.

If the fourth hypothesis were carried through in the same way as the three preceding ones,
it would give $X=0$, $Y=-0{,}000\,0168$, from which the following corrected values of $x$, $y$
would be obtained
\[
\begin{aligned}x &= \log P = 0{,}025\,6331 \\ &\quad\text{(the same as in the fourth hypothesis)}
\\
y &= \log Q = 0{,}750\,8973.
\end{aligned}
\]
\noindent If a fifth hypothesis were based on these values, the solution would attain the utmost
precision the tables allow, but the resulting elements would scarcely differ perceptibly
from those given by the fourth hypothesis.

\Needspace{19\baselineskip}
To obtain the complete elements, it only remains to compute the orientation of
the orbital plane. The rules of Art. 149 give
\begin{center}\begingroup\footnotesize\setlength{\tabcolsep}{3pt}\renewcommand{\arraystretch}{1.3}
\begin{tabular}{ll}\multicolumn{2}{c}{from the first position}\\
$g$\quad .\,.\,.\,.\,.\,.\,.\,.\,.\,.\,.\,.\,.\,.\,.\,.\,.\,.\,.\,.\,.\,.\,.\,.\,.\,.\,.\,. & $354^0\,9'\,44''{,}22$\\
$h$\quad .\,.\,.\,.\,.\,.\,.\,.\,.\,.\,.\,.\,.\,.\,.\,.\,.\,.\,.\,.\,.\,.\,.\,.\,.\,.\,.\,. & $261\quad56\quad6{,}94$\\
$i$\quad .\,.\,.\,.\,.\,.\,.\,.\,.\,.\,.\,.\,.\,.\,.\,.\,.\,.\,.\,.\,.\,.\,.\,.\,.\,.\,.\,. & $10\quad37\quad33{,}02$\\
$\AscendingNode$\quad .\,.\,.\,.\,.\,.\,.\,.\,.\,.\,.\,.\,.\,.\,.\,.\,.\,.\,.\,.\,.\,.\,.\,.\,.\,.\,. & $80\quad58\quad49{,}06$\\
Distance of perihelion from the ascending node & $65\quad2\quad4{,}47$\\
Longitude of perihelion\quad .\,.\,.\,.\,.\,.\,.\,.\,.\,.\,.\,.\,. & $146\quad0\quad53{,}53$\\
\noalign{\smallskip}\multicolumn{2}{c}{from the third position}\\
$g''$\quad .\,.\,.\,.\,. & $57^0\,5'\,0''{,}91$\\
$h''$\quad .\,.\,.\,.\,. & $161\quad0\quad1{,}61$\\
.\,.\,.\,.\,.\,.\,. & $10\quad37\quad33{,}00$\\
.\,.\,.\,.\,.\,.\,. & $80\quad58\quad49{,}10$\\
.\,.\,.\,.\,.\,.\,. & $65\quad2\quad4{,}52$\\
.\,.\,.\,.\,.\,.\,. & $146\quad0\quad53{,}62$.
\end{tabular}
\endgroup\end{center}
\noindent Taking the mean, we therefore set $i=10^{0}\,37'\,33''{,}01$, $\AscendingNode=80^{0}\,58'\,49''{,}08$,
longitude of perihelion $=146^{0}\,0'\,53''{,}57$. Finally, the mean longitude for the beginning
of 1806 will be $=108^{0}\,36'\,46''{,}08$.
